\chapter*{Preface}
\addcontentsline{toc}{chapter}{Preface}

Protein science, statistical physics, graphical modeling, and modern machine learning repeatedly
use the same words for different mathematical objects.  A coupling may mean a term in a joint
Hamiltonian, a conditional dependence parameter, a message map, or an experimentally measured
epistatic effect.  A graph may encode model factorization, computational communication, an
operator-valued response, or physical contact.  Softmax may normalize configurations, residue
identities, source positions, memories, or latent edge states.  The formulas can look nearly
identical while their scientific claims remain incompatible.

This book develops a typed mathematical language for navigating those translations.  Its central
object is a reference-centered operator-valued graph: every sequence position carries a categorical
contrast space, and every directed edge carries a map between such spaces.  This construction
places Potts couplings, protein-language-model responses, attention routes, latent memories,
spectral summaries, and biological interventions in common coordinates without declaring them
equivalent.

The purpose is not to attach physical language to every neural computation.  It is to identify
exactly what additional assumptions turn a score into an energy coordinate, a collection of
conditionals into a joint law, a route into an interaction, a scalar network into an operator
graph, and a predictive association into biological evidence.  Where those assumptions fail, the
book treats the failure as information rather than as a defect to conceal.

The material grew from a theory paper and retains theorem--proof precision, but its natural scope
is broader.  The translation dictionary, gauge geometry, static--dynamic decomposition,
higher-order hierarchy, graph probability, model correspondences, spectral analysis,
realizability constraints, and falsification program are parts of one framework.  Presenting them
as a book allows each layer to keep its proper meaning instead of forcing every distinction into
the narrow arc of a single research claim.

\section*{How to read this book}

Part I establishes the cross-disciplinary vocabulary and the hierarchy of claims.  Part II is the
theorem--proof core: it builds the operator geometry and develops static, dynamic, higher-order,
and probability-exact edge objects.  Part III has a mixed role.  It contains exact model
correspondences where their assumptions permit them, then uses the same coordinates as a
structured audit of Potts models, Gaussian models, Hopfield systems, RBMs, Transformers, and the
three generations of AlphaFold before turning to MSA estimators and graph spectra.  In particular,
the AlphaFold comparison is an architectural reading of published systems, not an equivalence
theorem about their algorithms.  Part IV asks which objects are realizable
in trainable architectures and what evidence is required for statistical, physical, or biological
interpretation.

Readers interested primarily in protein language models can begin with Chapters 1, 4, 7, 8, and 9.
Readers coming from statistical physics may prefer Chapters 2, 3, 6, 7, 8, and 10.  Readers designing
sparse or interpretable architectures should read Chapters 3, 4, 8, 11, and 12 before treating an
internal route as an identified interaction.

The front-matter object map gives a lighter first route through the notation.  On a first reading,
the opening prose, boxed equations, worked cases, and codas carry the conceptual argument.  Proofs
and mathematical interludes can be read on a second pass; they establish the exact boundaries of
the claims rather than introducing a separate storyline.
